\documentclass[12pt,a4paper]{article}

\usepackage{fontawesome5}
\providecommand{\faPencilAlt}{\faIcon{pencil-alt}}

\usepackage{styles/couleurs}
\usepackage{styles/mise_en_page}
\usepackage{styles/boites}
\usepackage{styles/maths}
\usepackage{styles/figures}

\usepackage{enumitem}
\usepackage{array}
\usepackage{tabularx}
\usepackage{amsmath}
\usepackage{amssymb}
\usepackage{tikz}
\usepackage{xparse}

\usetikzlibrary{calc,angles,quotes,arrows.meta}

\setlist[itemize]{leftmargin=1.4em,itemsep=0.15em,topsep=0.2em}
\setlist[enumerate]{leftmargin=*,itemsep=0.22em,topsep=0.25em}

\newcommand{\lignepointillee}{\noindent\dotfill\par\vspace{0.45em}}

\newcommand{\lignesreponse}[1]{%
\begin{tcolorbox}[
enhanced,colback=white,colframe=NeutreC!40,boxrule=0.5pt,arc=1.5mm,
left=3mm,right=3mm,top=2mm,bottom=1mm,before skip=2pt,after skip=2pt]
\foreach \i in {1,...,#1}{\lignepointillee}
\end{tcolorbox}}

\newcommand{\titreseq}[1]{%
\begin{center}
{\sffamily\LARGE\bfseries\color{TitreC}Séquence 6 -- Cosinus d'un angle aigu}\\[0.2em]
{\sffamily\large\bfseries\color{SavoirC}#1}
\end{center}}

\begin{document}

\titreseq{Activité 6 -- Choisir le bon triangle et la bonne méthode}

\section*{Partie A -- Préparer avant de calculer}

Pour chaque situation, complète le tableau avant d'effectuer les calculs.

\renewcommand{\arraystretch}{1.45}
\begin{tabularx}{\textwidth}{
|>{\centering\arraybackslash}p{1.1cm}
|>{\raggedright\arraybackslash}p{3cm}
|>{\raggedright\arraybackslash}p{2.6cm}
|>{\raggedright\arraybackslash}p{2.6cm}
|>{\centering\arraybackslash}p{2.2cm}
|X|}
\hline
 & \textbf{Triangle rectangle} & \textbf{Hypoténuse} &
\textbf{Côté adjacent} & \textbf{Longueur cherchée} &
\textbf{Multiplier ou diviser ?}\\
\hline
A & \(ABC\), rectangle en \(B\), \(AC=11\) cm,
\(\widehat{BAC}=42^\circ\), chercher \(AB\). & & & & \\
\hline
B & \(DEF\), rectangle en \(E\), \(DE=6{,}3\) cm,
\(\widehat{EDF}=51^\circ\), chercher \(DF\). & & & & \\
\hline
C & \(RST\), rectangle en \(S\), \(RT=14\) cm,
\(\widehat{SRT}=67^\circ\), chercher \(RS\). & & & & \\
\hline
D & \(JKL\), rectangle en \(K\), \(JK=9{,}1\) cm,
\(\widehat{KJL}=25^\circ\), chercher \(JL\). & & & & \\
\hline
\end{tabularx}

\section*{Partie B -- Calculer et rédiger}

Calcule les quatre longueurs. Arrondis au dixième.

\lignesreponse{12}

\section*{Partie C -- Une figure composée}

Le toit \(MNP\) est isocèle en \(M\).
La hauteur \([MH]\) coupe \([NP]\) en \(H\).

On donne :

\[
MN=7{,}8\text{ m}
\qquad\text{et}\qquad
\widehat{MNH}=34^\circ.
\]

\begin{center}
\begin{tikzpicture}[scale=0.85]
  \coordinate (N) at (0,0);
  \coordinate (P) at (8,0);
  \coordinate (H) at (4,0);
  \coordinate (M) at (4,4.3);

  \draw[line width=0.9pt] (N)--(M)--(P)--cycle;
  \draw[line width=0.9pt] (M)--(H);
  \draw (4,0.38)--(4.38,0.38)--(4.38,0);

  \node[left] at (N) {$N$};
  \node[right] at (P) {$P$};
  \node[above] at (M) {$M$};
  \node[below] at (H) {$H$};

  \node[above left] at (2,2.35) {$7{,}8$ m};
  \node[below] at (2,0) {$NH\,?$};

  \draw[ActionC,line width=1pt] (0.85,0) arc[start angle=0,end angle=34,radius=0.85];
  \node[ActionC] at (1.2,0.42) {$34^\circ$};
\end{tikzpicture}
\end{center}

\begin{enumerate}
  \item Dans quel triangle rectangle faut-il travailler ?
  \item Calcule \(NH\).
  \item Déduis la largeur totale \(NP\).
  \item Explique pourquoi on ne peut pas appliquer directement le cosinus
        dans le triangle \(MNP\).
\end{enumerate}

\lignesreponse{8}

\section*{Partie D -- Trois erreurs fréquentes}

Associe chaque erreur à une explication, puis corrige.

\begin{enumerate}
  \item \(\cos(\widehat A)=\dfrac{\text{côté opposé}}{\text{hypoténuse}}\).
  \item Pour trouver une hypoténuse, on multiplie le côté adjacent par le cosinus.
  \item On arrondit \(\cos(37^\circ)\) à \(0{,}8\) avant de poursuivre le calcul.
\end{enumerate}

Explications possibles :

\begin{itemize}
  \item le résultat risque d'être trop imprécis ;
  \item la définition du cosinus n'utilise pas le côté opposé ;
  \item le résultat obtenu serait plus petit que le côté adjacent.
\end{itemize}

\lignesreponse{6}

\end{document}
